\chapter{Виртуальные локальные сети}
\label{ch:chap2_vlan}
В данной лабораторной работе для организации межвиланной маршрутизации был выбран метод Router-on-a-Stick, основанный на использовании подинтерфейсов маршрутизатора. Такой подход обусловлен тем, что в лабораторной топологии применяются коммутаторы второго уровня (Layer 2), не поддерживающие маршрутизацию между VLAN. В связи с этим функции шлюза для каждой виртуальной локальной сети были перенесены на отдельный маршрутизатор, где для каждой VLAN был создан собственный подинтерфейс с соответствующим IP-адресом.

\section{Subinterfaces}
\begin{figure}[ht!]
    \centering
    \includegraphics[width=1\textwidth]{screenshots/2_1_0.jpg}
    \caption{Создание нужной топологии локальной сети для конфигурации VLAN}
    \label{fig:2_1_0}
\end{figure}
\clearpage
\begin{figure}[ht!]
    \centering
    \includegraphics[width=1\textwidth]{screenshots/2_1_1.jpg}
    \caption{Создание VLAN на коммутаторе Switch2}
    \label{fig:2_1_1}
\end{figure}
\clearpage
\begin{figure}[ht!]
    \centering
    \includegraphics[width=1\textwidth]{screenshots/2_1_2.jpg}
    \caption{Проверка, что VLAN с идентификаторами 10 и 20 созданы на коммутаторе Switch2}
    \label{fig:2_1_2}
\end{figure}
\clearpage
\begin{figure}[ht!]
    \centering
    \includegraphics[width=1\textwidth]{screenshots/2_1_3.jpg}
    \caption{Конфигурация портов на коммутаторе Switch2}
    \label{fig:2_1_3}
\end{figure}
\clearpage
\begin{figure}[ht!]
    \centering
    \includegraphics[width=1\textwidth]{screenshots/2_1_4.jpg}
    \caption{Проверка соответствия портов коммутатора Switch2 VLAN и trunk}
    \label{fig:2_1_4}
\end{figure}
\clearpage
\begin{figure}[ht!]
    \centering
    \includegraphics[width=1\textwidth]{screenshots/2_1_5.jpg}
    \caption{Такие же настройки только для коммутатора Switch3}
    \label{fig:2_1_5}
\end{figure}
\clearpage
\begin{figure}[ht!]
    \centering
    \includegraphics[width=1\textwidth]{screenshots/2_1_6.jpg}
    \caption{Настройки trunk на портах коммутатора Switch1}
    \label{fig:2_1_6}
\end{figure}
\clearpage
\begin{figure}[ht!]
    \centering
    \includegraphics[width=1\textwidth]{screenshots/2_1_7.jpg}
    \caption{Включение физического интерфейса на роутере}
    \label{fig:2_1_7}
\end{figure}
\clearpage
\begin{figure}[ht!]
    \centering
    \includegraphics[width=1\textwidth]{screenshots/2_1_8.jpg}
    \caption{Настройка подинтерфейсов VLAN на роутере}
    \label{fig:2_1_8}
\end{figure}
\clearpage
\begin{figure}[ht!]
    \centering
    \includegraphics[width=1\textwidth]{screenshots/2_1_9.jpg}
    \caption{Установка ip-адреса VPC с шлюзом по умолчанию}
    \label{fig:2_1_9}
\end{figure}
\clearpage
\begin{figure}[ht!]
    \centering
    \includegraphics[width=1\textwidth]{screenshots/2_1_10.jpg}
    \caption{Проверка доступности узлов VPC из одних/разных VLAN}
    \label{fig:2_1_10}
\end{figure}
\clearpage

\section{Выводы по главе}
В данной главе была рассмотрена конфигурация VLAN с использованием подинтерфейсов на маршрутизаторе. Были созданы VLAN с идентификаторами 10 и 20 на коммутаторах, настроены порты для соответствующих VLAN.
Между коммутаторами и маршрутизатором были настроены trunk-соединения по стандарту IEEE 802.1Q, обеспечивающие передачу трафика нескольких VLAN по одному физическому каналу, а также настроены подинтерфейсы для обеспечения маршрутизации между VLAN. В результате были проверены настройки и подтверждена доступность узлов VPC из разных VLAN, что демонстрирует успешную конфигурацию сети.
\endinput
